\documentclass[a4paper]{article}

\usepackage{graphicx}
\usepackage{amsmath,amssymb}
\usepackage{float}
\usepackage{natbib}

\title{Proportional controller}
\date{}
\author{Mohammad Rahmani}

\begin{document}
   \maketitle

   A proportional controller is a feedback controller whose control output is proportional to the current control error \cite{astrom-2008-feedback-systems-an-introduction-for-scientists-and-engineers}.

   \section{Formulation}
      Let the reference value be $r(t)$ and the measured system output be $y(t)$. The control error is
      \[
         e(t)=r(t)-y(t).
      \]

      A proportional controller generates the control input
      \[
         u(t)=K_P e(t),
      \]
      where $K_P$ is the proportional gain. Therefore, a larger error produces a proportionally larger control action.

   \subsection{Vector form}
      For a multi-dimensional state,
      \[
         \mathbf{e}(t)=\mathbf{r}(t)-\mathbf{y}(t),
         \qquad
         \mathbf{u}(t)=K_P\mathbf{e}(t),
      \]
      where $K_P$ may be a scalar gain or a gain matrix.

   \subsection{Position-control loop}
      For position control, the position error can be converted into a velocity setpoint:
      \[
         \mathbf{e}_p
         =
         \mathbf{p}_{\mathrm{sp}}
         -
         \mathbf{p},
         \qquad
         \mathbf{v}_{\mathrm{sp}}
         =
         K_P\mathbf{e}_p.
      \]
      Thus, the farther the system is from the desired position, the larger the commanded velocity becomes. As the position error approaches zero, the commanded velocity also approaches zero.

   \subsection{Closed-loop form}
      For a plant with transfer function $G(s)$ and proportional gain $K_P$, the closed-loop transfer function under unity negative feedback is
      \[
         \frac{Y(s)}{R(s)}
         =
         \frac{K_P G(s)}
              {1+K_P G(s)}.
      \]

      Increasing $K_P$ generally increases the strength of the corrective response, but an excessively large gain can increase overshoot, oscillation, or instability depending on the plant dynamics \cite{astrom-2008-feedback-systems-an-introduction-for-scientists-and-engineers}.

   \bibliographystyle{plainnat}
   \bibliography{/home/donkarlo/Dropbox/repo/nd_shared_research/bibliography/bibliography.bib}
\end{document}
